\section{Threats to validity}
\noindent \textbf{External Validity}.
One of the major threats to the external validity of our work is the data collection internal. We collect data in an 11 years old period from 2011 to 2022 to coverall all possible studies published during this period. Another source of threat to external validity is the coverage of the input data types in software vulnerability detection. To tackle this threat, we focused on source code snippets as well as repository data, i.e., data that can be extracted from open source including GitHub and CVE. 

\noindent \textbf{Internal Validity}. 
One of the major threats to the internal validity of our work is the automatic collection of data for analysis. Our technique for data collection is automated at its initial steps in which we may miss some important vulnerability detection papers. In fact, we developed a set of scripts that allowed us to extract papers given an 11-year period. Even though the subsequent steps are manually supervised, still the automatic filtering suffers from this issue. The second important issue with our data collection is possible bias. The root cause of bias is unavoidable disagreements during paper classification as two researchers have worked together to categorize papers based on title, abstract, and content of papers. In case the researchers could not come up with an agreement, the third researcher join to resolve the differences which somehow relaxes the issue. 

\noindent \textbf{Construct Validity}. The major threat to the construct validity of our survey is the level of granularity of the analysis we conducted for each primary study. For each study, we deeply analyzed the artifacts explained in the paper and show their distribution as tables and graphs. For example, in terms of sources for benchmark data, we deeply analyzed 17 sources accounting for 53 primary studies overall. The second threat to the construct validity is the degree of coverage for each primary study. We analyzed each primary study from 5 aspects including input data, input representation, embedding techniques, models, vulnerability types, and whether the study supports the interpretability of vulnerability detection models or not.  

%\song{check other surveys to add more limitations here} \nima{done}
% Study Selection Bias. Another threat might be the possibility of missing DL related studies during the literature search and selection phase. We are unable to identify and retrieve all relevant publications considering the many publication venues publishing SE relevant papers. Therefore, we carefully selected 21 publication venues, including conference proceedings, symposiums, and journals, which can cover many prior works in SE. Besides, we identified our search terms and formed a search string by combining and amending the search strings used in other literature reviews on DL. These could keep the number of missed primary studies as small as possible.

